\chapter*{Introduction générale}
\addcontentsline{toc}{chapter}{Introduction générale}
\markboth{Introduction générale}{}

Dans un monde en constante évolution technologique, où les smartphones sont devenus des outils indispensables du quotidien, le secteur de la santé connaît une transformation numérique profonde. La digitalisation des services de santé, communément appelée e-santé, ouvre de nouvelles perspectives pour améliorer la qualité des soins et faciliter le suivi médical des patients.

L'un des défis majeurs auxquels sont confrontés les patients, particulièrement ceux souffrant de maladies chroniques ou suivant des traitements de longue durée, est la gestion efficace de leurs médicaments. L'observance thérapeutique, c'est-à-dire le respect par le patient des prescriptions médicales, reste un problème crucial dans le système de santé moderne. Selon l'Organisation Mondiale de la Santé, environ 50\% des patients atteints de maladies chroniques ne suivent pas correctement leur traitement, ce qui entraîne une diminution de l'efficacité thérapeutique et une augmentation des coûts de santé.

Les difficultés rencontrées par les patients sont multiples : oubli de la prise des médicaments, confusion entre différents traitements, gestion des dates de péremption, suivi des quantités disponibles, et coordination avec les membres de la famille ou les aidants. Ces problématiques sont encore plus complexes pour les personnes âgées, les patients polymorbides, ou les parents gérant les traitements de leurs enfants.

Face à ces enjeux, les solutions technologiques mobiles représentent une opportunité significative pour améliorer la gestion des médicaments au quotidien. C'est dans ce contexte que s'inscrit notre projet de fin d'études, qui consiste à développer \textbf{AFYA}, une application mobile de gestion des médicaments destinée à faciliter le suivi des traitements médicaux et à améliorer l'observance thérapeutique.

AFYA est une solution complète qui permet aux utilisateurs de gérer leurs prescriptions, de recevoir des rappels personnalisés pour la prise de médicaments, de suivre l'historique de leurs traitements, et de gérer efficacement leur boîte à pharmacie. L'application intègre également des fonctionnalités collaboratives permettant aux membres d'une famille ou aux professionnels de santé de superviser les traitements des patients, favorisant ainsi une meilleure coordination des soins.

Le projet a été réalisé au sein de la société \textbf{IDVEY}, spécialisée dans le développement de solutions digitales innovantes. Sur le plan technique, AFYA a été développée en utilisant des technologies modernes et performantes : \textbf{Flutter} pour le développement de l'application mobile multi-plateforme (Android et iOS), et \textbf{Spring Boot} pour le développement du backend REST API, garantissant ainsi une architecture robuste, sécurisée et scalable.

Le développement du projet s'est étalé sur une période de plusieurs mois, structurée selon la méthodologie \textbf{Agile Scrum} en 6 sprints répartis sur 3 releases. Cette approche itérative et incrémentale nous a permis de livrer progressivement des fonctionnalités à valeur ajoutée, tout en maintenant une grande flexibilité face aux changements et aux retours utilisateurs.

Ce rapport retrace les différentes étapes de conception et de réalisation du projet AFYA. Il s'articule autour de cinq chapitres principaux :

\begin{itemize}
    \item \textbf{Le premier chapitre} présente le contexte général du projet. Nous y décrivons l'organisme d'accueil IDVEY, la problématique abordée, l'étude de l'existant et la critique des solutions actuelles. Nous présentons ensuite notre solution proposée ainsi que la méthodologie de travail adoptée pour la réalisation du projet.

    \item \textbf{Le deuxième chapitre} est consacré à l'analyse et à la spécification des besoins. Nous y identifions les différents acteurs du système, définissons les besoins fonctionnels et non fonctionnels, et présentons le diagramme de cas d'utilisation global. Nous exposons également le backlog produit complet, la répartition des sprints en releases, ainsi que l'architecture logicielle et l'environnement technique choisis.

    \item \textbf{Le troisième chapitre} détaille la première release du projet, centrée sur la mise en place des fonctionnalités de base. Nous y présentons les travaux réalisés durant les sprints 1 et 2, incluant l'authentification, la gestion des groupes, l'ajout de médicaments, la gestion des traitements et des prescriptions, ainsi que la gestion des boîtes à pharmacie. Cette release constitue le socle fondamental de l'application.

    \item \textbf{Le quatrième chapitre} présente la deuxième release, axée sur la gestion collaborative et les notifications. Les sprints 3 et 4 ont permis d'implémenter les fonctionnalités de suivi des doses, de marquage des prises de médicaments par les responsables et gestionnaires, ainsi qu'un système complet de notifications et de rappels pour améliorer l'observance thérapeutique.

    \item \textbf{Le cinquième chapitre} expose la troisième et dernière release, dédiée aux fonctionnalités avancées et à l'enrichissement de l'expérience utilisateur. Les sprints 5 et 6 ont introduit la gestion complète du profil utilisateur, l'historique des traitements, les opérations CRUD avancées sur les médicaments, prescriptions et traitements, ainsi que la gestion fine des boîtes à pharmacie.
\end{itemize}

Enfin, une conclusion générale vient clôturer ce rapport en dressant un bilan du travail accompli, en soulignant les compétences acquises durant ce projet, et en proposant des perspectives d'évolution pour l'application AFYA.

\vspace{0.5cm}
